\section*{Chapter \ref{ch:m2:loans-clos-and-securitisation} --- Loans, CLOs and Securitisation}

\begin{solution}{exo:m2:loans-clos-and-securitisation:1}
$500/365 = 1.370$, $500/395 = 1.266$, $500/425 = 1.176$ and $500/450 = 1.111$.
\end{solution}

\begin{solution}{exo:m2:loans-clos-and-securitisation:2}
The loan pays a floating rate plus a fixed spread, reset each period from a short-term
rate that follows the central bank's policy rate; the bond's coupon was fixed at issue.
\end{solution}

\begin{solution}{exo:m2:loans-clos-and-securitisation:3}
The maintenance \omterm{def:m2:corporate-bonds:covenant}{covenants}: the right to act, to renegotiate or to accelerate, when the
borrower's ratios deteriorate, before it misses a payment. Lenders learn of trouble
later and have less say until a default.
\end{solution}

\begin{solution}{exo:m2:loans-clos-and-securitisation:4}
The par may fall by $500 - 1.18 \times 395 = 33.9$ million; at a loss of 30\% per
default that is USD~113 million of defaults, 22.6\% of the portfolio.
\end{solution}

\begin{solution}{exo:m2:loans-clos-and-securitisation:5}
The rule applies to securitisers, who transfer assets into a \omterm{def:m2:loans-clos-and-securitisation:sec}{securitisation}; an
open-market CLO manager neither originates the loans nor holds them before the vehicle
buys them, so it has nothing to retain in that sense. Risk retention addresses
originate-to-distribute: a lender that sells everything it makes has little reason to
lend carefully.
\end{solution}

\begin{solution}{exo:m2:loans-clos-and-securitisation:6}
The loans earn 7.5\% on USD~500 million; the notes cost 5.81\% on average on USD~450
million, and the fee 0.45\% on the whole. The equity, USD~50 million, collects the
difference on the whole pool: $37.5 - 2.25 - 26.14 = 9.11$ million, 18.2\% of its
investment, because its own small stake earns the spread earned on ten times as much.
\end{solution}

\begin{solution}{exo:m2:loans-clos-and-securitisation:7}
5.35\% a year instead of 4.78\%. After the reinvestment period, recoveries repay the
AAA notes instead of buying new loans, which lowers the notes faster than the par and
raises the coverage ratios, so the tests fail later.
\end{solution}

\begin{solution}{exo:m2:loans-clos-and-securitisation:8}
The CLO's AAA notes are rated on the structure: their safety comes from subordination
and tests against a pool of speculative-grade loans, and it depends on how defaults in
the pool are correlated, which a single company's bond does not. Their price can fall
far more in a credit crisis, when correlation and downgrades rise together, and they
are less liquid.
\end{solution}

\begin{solution}{pb:m2:loans-clos-and-securitisation:1}
\textbf{1.} USD~37.5 million; the fee USD~2.25 million, the notes USD~26.14 million.
\textbf{2.} USD~9.11 million, 18.2\%.
\textbf{3.} 19.5\% a year.
\textbf{4.} USD~27.5 million ($500 - 1.05 \times 450$).
\textbf{5.} USD~91.7 million, 18.3\% of the portfolio.
\textbf{6.} 13.1\%, 5.0\% and $-5.8\%$.
\textbf{7.} USD~2.0, 12.5 and 28.4 million.
\textbf{8.} 4.78\% a year.
\textbf{9.} In quarter 19, when the portfolio's par has fallen by USD~33.6 million.
\textbf{10.} USD~6.5 million of its USD~25 million is not repaid, and the equity receives
nothing from quarter 8.
\textbf{11.} It raises the cut-off to 5.35\%: recoveries then repay the notes.
\textbf{12.} Every failure turns junior cash into repayments of the AAA notes, which
lowers their balance faster than the pool loses par.
\textbf{13.} Each default lowers the par by the loss, 30\% of the loan against 60\% of
the bond, so it uses up half as much of the cushion.
\textbf{14.} Sell weak loans before they default, or buy loans below par, which the
tests may count at par; the first crystallises losses, the second adds risk, and
both are limited by the documents.
\textbf{15.} It returns the equity's share of the principal; if loans must be sold
below par, as in a stressed market, that final payment shrinks or disappears.
\textbf{16.} With the realised default rates on \omterm{def:m2:loans-clos-and-securitisation:loan}{leveraged loans}, including distressed
exchanges, in past recessions; defaults cluster in a recession, so a few years well
above the average do more damage than the same total spread evenly.
\textbf{17.} The equity: investors who can bear large losses and value the excess
spread; the AAA notes: investors who need high ratings and floating-rate assets, such
as banks and insurers.
\textbf{18.} Buy protection on a high-yield \omterm{def:m2:credit-indices-and-tranches:index}{credit index} or on its junior tranche, which
lose in the same default wave; the hedge is imperfect because the CLO's loans are not
the index's names.
\textbf{19.} Named result: \emph{the test that trips} cuts the equity off at a
constant default rate of 4.78\% a year, first in quarter 19, when the par has fallen
by USD~33.6 million.
\textbf{20.} It takes cash from the equity and gives it to the senior notes as soon as
losses eat into the cushion.
\end{solution}

\begin{solution}{iq:m2:loans-clos-and-securitisation:1}
A vehicle holding an actively managed pool of \omterm{def:m2:loans-clos-and-securitisation:loan}{leveraged loans}, funded by notes rated
from AAA to BB and by unrated equity; the cash flows pay the notes in order and the
rest to the equity. Banks, insurers, funds and pension funds hold the rated notes; the
equity goes to investors who can bear first losses.
\iqlookfor{SPV, manager, tranches, investors by risk.}
\end{solution}

\begin{solution}{iq:m2:loans-clos-and-securitisation:2}
On each payment date, the ratio of loan par to the notes down to a class, and of
interest collected to interest due, are compared with triggers; if one fails, the cash
that would have gone further down repays the senior notes until the ratio is restored.
Defaults, losses and downgraded loans move the ratios.
\iqlookfor{par and interest ratios, triggers, diversion.}
\end{solution}

\begin{solution}{iq:m2:loans-clos-and-securitisation:3}
Model the pool's defaults, recoveries, prepayments and rates, by scenario or
simulation, with loan-level detail if available; run the waterfall for each path; the
equity's value is the discounted residual cash, at a rate reflecting its risk, and
its sensitivity to the default and recovery assumptions is the main risk. Check the
result against secondary prices.
\iqlookfor{scenarios through the actual waterfall.}
\end{solution}

\begin{solution}{iq:m2:loans-clos-and-securitisation:4}
An index tranche is synthetic, static, of fixed maturity and marked daily from index
spreads; a CLO tranche is cash, backed by a managed, changing pool of loans, with
reinvestment and tests that move cash between classes, and trades rarely. Its risk
depends on the manager and the documents, not only on default correlation.
\iqlookfor{static versus managed, tests, liquidity.}
\end{solution}

\begin{solution}{iq:m2:loans-clos-and-securitisation:5}
Rising leverage and weaker \omterm{def:m2:corporate-bonds:covenant}{covenants} among borrowers, floating-rate debt that makes
them vulnerable to higher rates, concentration in some sectors, and the chain of
holders (CLOs, funds, insurers and banks) through which losses spread; funds that
promise liquidity against illiquid loans can face runs.
\iqlookfor{borrower fragility and where the losses land.}
\end{solution}

\begin{solution}{iq:m2:loans-clos-and-securitisation:6}
Encode each deal's waterfall as data (classes, tests, triggers, fees, dates) run by one
engine rather than per-deal code; vectorise over scenarios; share the loan-level cash
flows across deals that hold the same loans; run in parallel across deals and nodes;
validate each deal against its trustee reports.
\iqlookfor{waterfall as data, reuse, parallelism, validation.}
\end{solution}
